\documentclass[11pt,a4paper]{article}
\usepackage[T1]{fontenc}
\usepackage{lmodern}
\usepackage[margin=27mm]{geometry}
\usepackage{amsmath,amssymb,amsthm,mathtools}
\usepackage{microtype}
\usepackage{booktabs,array}
\usepackage{enumitem}
\usepackage[hidelinks]{hyperref}
\setlength{\parskip}{5pt plus 1pt}
\setlength{\parindent}{0pt}
\setlength{\emergencystretch}{3em}
\setlist{itemsep=2pt,topsep=3pt}
\raggedbottom
\newtheorem{theorem}{Theorem}[section]
\newtheorem{lemma}[theorem]{Lemma}
\newtheorem{corollary}[theorem]{Corollary}
\newtheorem{proposition}[theorem]{Proposition}
\theoremstyle{definition}
\newtheorem{definition}[theorem]{Definition}
\newtheorem{example}[theorem]{Example}
\theoremstyle{remark}
\newtheorem{remark}[theorem]{Remark}
\newcommand{\Z}{\mathbb{Z}}
\newcommand{\N}{\mathbb{N}}
\newcommand{\Q}{\mathbb{Q}}
\newcommand{\F}{\mathbb{F}}
\newcommand{\ps}[1]{#1[[x]]}
\newcommand{\coeff}[2]{[x^{#1}]\,#2}
\newcommand{\hl}[1]{\textbf{#1}}
\newcommand{\para}[1]{\par\medskip\noindent\textbf{#1}\enspace}
\hypersetup{pdftitle={Congruences modulo 4 for the series of OEIS A389537},pdfauthor={Dev D. Goyal},pdfsubject={OEIS A389537}}
\title{Congruences modulo $4$ for the series with $A\bigl(A(x)-x\bigr)=A(x)^2$\\[6pt]\large A proof of a conjecture in OEIS A389537}
\author{Dev D. Goyal}
\date{October 2026}
\begin{document}
\maketitle
\begin{abstract}
The sequence A389537 of the On-Line Encyclopedia of Integer Sequences (OEIS) lists the coefficients of the power series
$A(x)=x+x^2+2x^3+4x^4+8x^5+14x^6+\cdots$ that satisfies $A\bigl(A(x)-x\bigr)=A(x)^2$. Paul~D.~Hanna conjectured that,
for $n>2$, the coefficient $a(n)$ leaves remainder $2$ on division by $4$ when $n$ is $3$ times a power of $2$, and is
divisible by $4$ otherwise. We prove this. As the entry records, $A$ undoes the substitution of the sparse series
$R(x)=x-x^2+x^4-x^8+\cdots$. Modulo $2$ this gives $A\equiv x+x^2$. Writing $A=x+x^2+2H$ and working modulo $4$
leads to the equation $H+H^2=x^3$ modulo $2$, whose solution is $x^3+x^6+x^{12}+x^{24}+\cdots$. The results have been
checked by computer and proved formally in the Lean proof assistant.
\end{abstract}
\begin{center}\begin{minipage}{0.88\linewidth}\small
\textbf{How this paper was produced.} This paper reports the result of an AI research effort directed by
Dev~D.~Goyal. ChatGPT (OpenAI), using the GPT-6 model with Pro reasoning and working autonomously from prompts by
Dev~D.~Goyal, found the proof, wrote its Lean~4 formalization, and wrote a first draft of the paper.
Dev~D.~Goyal chose the problems and asked for proofs, formal verification and revisions, but did not supply proof
ideas. Claude Opus~5.5, used with Claude Code, then checked the work independently: it rebuilt the formal proofs from
source and confirmed that they rely only on Lean's three standard axioms, recomputed the sequence from its definition
compared it with the published terms and tested every claim numerically, and read every step of the proofs. It also wrote the present
version of this paper.
Dev~D.~Goyal has reviewed the proofs and is responsible for them; this review is not an independent expert validation.
The paper has not been peer reviewed.
The section ``About this paper'' at the end gives the details.
\end{minipage}\end{center}

\section{Introduction}

\subsection{The sequence}\label{sec:sequence}
The \emph{On-Line Encyclopedia of Integer Sequences} (OEIS)~\cite{oeisdb} is a large public database of integer
sequences. Many
entries also record \emph{conjectures}: patterns that the contributor observed in the data but did not prove.

This paper is about entry A389537~\cite{oeis}, contributed by Paul~D.~Hanna in November 2025. It defines a power
series $A(x)=a(1)\,x+a(2)\,x^2+a(3)\,x^3+\cdots$ with $a(1)=1$ by the equation
\begin{equation}\label{eq:oeis}
A\bigl(A(x)-x\bigr)=A(x)^2 .
\end{equation}
The left side is the series obtained by substituting $A(x)-x=x^2+\cdots$ for $x$ in $A$
(Section~\ref{sec:background} explains substitution). The right side is an ordinary square. Theorem~\ref{thm:exist} shows that there is exactly one such series.
Its first coefficients, starting at $n=1$, are 1, 1, 2, 4, 8, 14, 16, $-$20, $-$232, $-$1100, $-$4120, $-$13494, \dots{} Some later terms
are negative. The entry reads, in part:
\begin{quote}\small\ttfamily\raggedright G.f. A(x) satisfies A( A(x) - x ) = A(x)\^{}2.\par\smallskip Conjecture: for n > 2, a(n) == 2 (mod 4) when n = 3*2\^{}k for k >= 0 otherwise a(n) is divisible by 4.\par\smallskip For n > 7, it appears that the run-length of signs is either 9 or 10 (with no obvious pattern).\par\smallskip (5) A(R(x)) = x where R(x) = x - R(x\^{}2) = Sum\_\{n>=0\} (-1)\^{}n * x\^{}(2\^{}n).\par\smallskip\hfill\textrm{--- OEIS A389537, Paul D. Hanna, November 2025}\end{quote}

Here ``G.f.''\ stands for \emph{generating function}, the series whose coefficients are the terms of the sequence.
To avoid confusion with the equation numbers of this paper, we refer to the formula numbered (5) in the entry as (E5).
Formula (E5) says that $A$ undoes the substitution of the series $R(x)=x-x^2+x^4-x^8+x^{16}-\cdots$.

\subsection{The question}
The conjecture concerns the remainders of $a(n)$ on division by $4$. It says that, from $n=3$ on, $a(n)$ is divisible
by $4$, except when $n$ is one of $3,6,12,24,48,\dots$ (three times a power of $2$), where $a(n)$ leaves remainder $2$.
In binary, these exceptional $n$ are exactly the numbers written as $11$ followed by zeros. The table shows the
remainders of the first $48$ terms; the exceptional indices are printed in bold.
\begin{center}\small\setlength{\tabcolsep}{3.2pt}
\begin{tabular}{>{\raggedright\arraybackslash}p{6.4em}>{\centering\arraybackslash}p{1.15em}>{\centering\arraybackslash}p{1.15em}>{\centering\arraybackslash}p{1.15em}>{\centering\arraybackslash}p{1.15em}>{\centering\arraybackslash}p{1.15em}>{\centering\arraybackslash}p{1.15em}>{\centering\arraybackslash}p{1.15em}>{\centering\arraybackslash}p{1.15em}>{\centering\arraybackslash}p{1.15em}>{\centering\arraybackslash}p{1.15em}>{\centering\arraybackslash}p{1.15em}>{\centering\arraybackslash}p{1.15em}>{\centering\arraybackslash}p{1.15em}>{\centering\arraybackslash}p{1.15em}>{\centering\arraybackslash}p{1.15em}>{\centering\arraybackslash}p{1.15em}}
\toprule
$n$ & 1 & 2 & \hl{3} & 4 & 5 & \hl{6} & 7 & 8 & 9 & 10 & 11 & \hl{12} & 13 & 14 & 15 & 16 \\ \midrule
$a(n)\bmod 4$ & 1 & 1 & \hl{2} & 0 & 0 & \hl{2} & 0 & 0 & 0 & 0 & 0 & \hl{2} & 0 & 0 & 0 & 0 \\
\bottomrule
\end{tabular}
\\[4pt]
\begin{tabular}{>{\raggedright\arraybackslash}p{6.4em}>{\centering\arraybackslash}p{1.15em}>{\centering\arraybackslash}p{1.15em}>{\centering\arraybackslash}p{1.15em}>{\centering\arraybackslash}p{1.15em}>{\centering\arraybackslash}p{1.15em}>{\centering\arraybackslash}p{1.15em}>{\centering\arraybackslash}p{1.15em}>{\centering\arraybackslash}p{1.15em}>{\centering\arraybackslash}p{1.15em}>{\centering\arraybackslash}p{1.15em}>{\centering\arraybackslash}p{1.15em}>{\centering\arraybackslash}p{1.15em}>{\centering\arraybackslash}p{1.15em}>{\centering\arraybackslash}p{1.15em}>{\centering\arraybackslash}p{1.15em}>{\centering\arraybackslash}p{1.15em}}
\toprule
$n$ & 17 & 18 & 19 & 20 & 21 & 22 & 23 & \hl{24} & 25 & 26 & 27 & 28 & 29 & 30 & 31 & 32 \\ \midrule
$a(n)\bmod 4$ & 0 & 0 & 0 & 0 & 0 & 0 & 0 & \hl{2} & 0 & 0 & 0 & 0 & 0 & 0 & 0 & 0 \\
\bottomrule
\end{tabular}
\\[4pt]
\begin{tabular}{>{\raggedright\arraybackslash}p{6.4em}>{\centering\arraybackslash}p{1.15em}>{\centering\arraybackslash}p{1.15em}>{\centering\arraybackslash}p{1.15em}>{\centering\arraybackslash}p{1.15em}>{\centering\arraybackslash}p{1.15em}>{\centering\arraybackslash}p{1.15em}>{\centering\arraybackslash}p{1.15em}>{\centering\arraybackslash}p{1.15em}>{\centering\arraybackslash}p{1.15em}>{\centering\arraybackslash}p{1.15em}>{\centering\arraybackslash}p{1.15em}>{\centering\arraybackslash}p{1.15em}>{\centering\arraybackslash}p{1.15em}>{\centering\arraybackslash}p{1.15em}>{\centering\arraybackslash}p{1.15em}>{\centering\arraybackslash}p{1.15em}}
\toprule
$n$ & 33 & 34 & 35 & 36 & 37 & 38 & 39 & 40 & 41 & 42 & 43 & 44 & 45 & 46 & 47 & \hl{48} \\ \midrule
$a(n)\bmod 4$ & 0 & 0 & 0 & 0 & 0 & 0 & 0 & 0 & 0 & 0 & 0 & 0 & 0 & 0 & 0 & \hl{2} \\
\bottomrule
\end{tabular}
\end{center}

\subsection{What this paper proves}
\textbf{Main result} (Theorem~\ref{thm:main}). \emph{The conjecture is true: for every $n>2$,
$a(n)\equiv2\pmod 4$ if $n=3\cdot2^k$ for some $k\ge0$, and $a(n)\equiv0\pmod 4$ otherwise.}

The key steps are as follows. The series $R$ satisfies the equation $R(x)+R(x^2)=x$. Modulo $2$, squaring a
series is the same as replacing $x$ by $x^2$, so this becomes $R+R^2\equiv x$, which says that $x+x^2$ undoes the
substitution of $R$. Hence $A\equiv x+x^2\pmod 2$: all coefficients from $x^3$ on are even. Writing them as
$A=x+x^2+2H$, we then compute modulo $4$. Equation \eqref{eq:oeis} reduces to the equation $H+H^2\equiv x^3$
modulo $2$, whose solution is $H\equiv x^3+x^6+x^{12}+x^{24}+\cdots$. A coefficient $a(n)=2\,\coeff{n}{H}$ with $n>2$
then leaves remainder $2$ on division by $4$ exactly when $\coeff{n}{H}$ is odd.

\subsection{What was known}
The entry records formula (E5), which identifies $A$ as the inverse of $R$; we prove it here (Theorem~\ref{thm:exist})
because the argument relies on it. We are not aware of an earlier proof of the conjecture. The entry's second
observation, about runs of equal signs, concerns the signs of the coefficients rather than their remainders, and we
do not address it.

\subsection{Motivation}
Series such as $R(x)=x-x^2+x^4-x^8+\cdots$, whose nonzero terms become ever sparser, are called \emph{lacunary}. Because
$R$ is related to $R(x^2)$ by a linear equation, its coefficients are tied to the binary digits of $n$. Patterns of this
kind are typical of \emph{automatic sequences}, sequences that can be generated by a finite machine reading the binary
digits of $n$~\cite{allouche}. A classical theorem of Christol~\cite{christol} explains when to expect them: a power series
with coefficients modulo $2$ satisfies a polynomial equation, whose coefficients are themselves polynomials in $x$,
exactly when its coefficients form such a sequence. Here
the reduction of $H$ modulo $2$ satisfies the polynomial equation $H+H^2=x^3$, and the answer indeed depends only on the
binary digits of $n$. The proof shows how such structure can be found in a sequence defined by a much less transparent
equation.

\subsection{How to read this paper}
Section~\ref{sec:background} explains the notions we need: power series, substitution, compositional inverses, and
remainders. Section~\ref{sec:results} states the results, and Section~\ref{sec:proofs} proves them.
Section~\ref{sec:checks} describes how the results were checked, by computer and by a formal proof in Lean.
Section~\ref{sec:discussion} discusses the scope of the results and open questions.
\section{Background}\label{sec:background}
This section collects the definitions and facts used later. Readers who know formal power series and congruences can skip it and return when needed.
\subsection{Formal power series}
A \emph{formal power series} is an expression
\[
U(x)=u_0+u_1x+u_2x^2+u_3x^3+\cdots ,
\]
determined by its sequence of coefficients $u_0,u_1,u_2,\dots$; see~\cite{niven,stanley}. No number is ever
substituted for $x$, so questions of convergence do not arise. We write $\coeff{n}{U}=u_n$ for the coefficient of $x^n$ in $U$; two series are equal
when all their coefficients are equal.

Series are added coefficient by coefficient and multiplied like polynomials:
\[
\coeff{n}{(UV)}=u_0v_n+u_1v_{n-1}+\cdots+u_nv_0 .
\]
Each coefficient of a product is a finite sum. For example, $(1-x)(1+x+x^2+x^3+\cdots)=1$, so $1+x+x^2+\cdots$ is
the \emph{multiplicative inverse} of $1-x$, written $\frac{1}{1-x}$. More generally, every series $D$ with constant term
$1$ has a unique multiplicative inverse $\frac1D$, whose coefficients are determined one at a time by $D\cdot\frac1D=1$
using only addition, subtraction and multiplication. Whenever a quotient of power series appears in this paper, its
denominator has constant term $1$, so the quotient is again a power series.

A \emph{commutative ring} is a set with an addition, a subtraction and a multiplication that satisfy the usual rules
of arithmetic; examples are $\Z$, $\Q$ and the integers modulo $m$. For a commutative ring $R$ we write $\ps{R}$ for
the set of power series with coefficients in $R$; for instance, writing $\F_2$ for the integers modulo $2$, $\ps{\F_2}$
consists of the series with coefficients modulo $2$. We write
$U(0)$ for the constant term $u_0$ of $U$, and $U=O(x^d)$ when the coefficients of $1,x,\dots,x^{d-1}$ in $U$ are all
zero. Thus $U-V=O(x^d)$ means that $U$ and $V$ have the same coefficient of $x^n$ for every $n<d$. A series with
constant term $0$ and coefficient of $x$ equal to $1$, that is, a series of the form $x+u_2x^2+u_3x^3+\cdots$, is
called \emph{normalized}.

\subsection{Substituting one series into another}
If $V(x)$ has constant term $0$, we can substitute $V$ for $x$ in $U$:
\[
U(V(x))=u_0+u_1V(x)+u_2V(x)^2+u_3V(x)^3+\cdots .
\]
This is called \emph{composition}, also written $U\circ V$. It makes sense because $V(x)^k$ has no terms below $x^k$,
so only the finitely many terms with $k\le n$ affect the coefficient of $x^n$. For example, if $U=x+x^2$ and $V=x^2$, then
$U(V(x))=x^2+x^4$. Composition is associative: $U\circ(V\circ W)=(U\circ V)\circ W$.

Substituting a series into itself repeatedly gives its \emph{iterates}: $U^{\circ2}(x)=U(U(x))$,
$U^{\circ3}(x)=U(U(U(x)))$, and so on, with $U^{\circ0}(x)=x$ and $U^{\circ1}=U$. They must not be confused with
\emph{powers} $U^2=U\cdot U$, $U^3=U\cdot U\cdot U$. (The OEIS sometimes writes $A^k(x)$ for the $k$-th iterate. In
this paper the small circle always means an iterate, and a plain exponent always means a power.)

\subsection{Undoing a substitution}
A normalized series can always be ``undone'' by another normalized series.

\begin{lemma}[Compositional inverse]\label{lem:inverse}
For every normalized series $U$ there is exactly one normalized series $V$ with $U(V(x))=x$, and this $V$ also
satisfies $V(U(x))=x$. If $U$ has integer coefficients, so does $V$.
\end{lemma}

We call $V$ the \emph{compositional inverse} of $U$. The OEIS calls it the \emph{series reversion} of $U$. For example,
the compositional inverse of $x-x^2$ is $C(x)=x+x^2+2x^3+5x^4+14x^5+\cdots$, whose coefficients are the Catalan numbers.
Indeed $C-C^2=x$, which says exactly that substituting $C$ into $x-x^2$ gives $x$.

\begin{proof}
Write $V=x+v_2x^2+v_3x^3+\cdots$ with unknown coefficients. In
$U(V)=V+u_2V^2+u_3V^3+\cdots$, the coefficient of $x^n$ (for $n\ge2$) equals $v_n$ plus terms that involve only
$v_2,\dots,v_{n-1}$. Setting it equal to $0$ therefore determines $v_n$ from the earlier coefficients, using only
additions, subtractions and multiplications. So there is exactly one $V$ with $U\circ V=x$, and it has integer
coefficients if $U$ does. Applying the same construction to $V$ gives a normalized $W$ with $V\circ W=x$. Then
$U=U\circ(V\circ W)=(U\circ V)\circ W=W$, so $V\circ U=V\circ W=x$ as well.
\end{proof}

\subsection{Remainders and congruences}
For an integer $m\ge2$ we write $a\equiv b\pmod m$ (``$a$ is congruent to $b$ modulo $m$'') when $a-b$ is divisible by
$m$, that is, when $a$ and $b$ leave the same remainder after division by $m$. The remainder of $a$ on division by
$m$ is called its \emph{residue} modulo $m$; it is one of $0,1,\dots,m-1$.

Addition and multiplication are compatible with congruences, so the residues $0,1,\dots,m-1$, added and multiplied
modulo $m$, form a ring, written $\Z/m\Z$. When $m=p$ is a prime we also write $\F_p$. In $\F_2=\{0,1\}$ we have $1+1=0$,
so addition and subtraction coincide.

For power series with integer coefficients, $U\equiv V\pmod m$ means that \emph{every} coefficient of $U-V$ is
divisible by $m$. Equivalently, the two series become equal when every coefficient is replaced by its residue; we call
this \emph{reducing} the series modulo $m$. The key fact is the following.

\emph{Reduction modulo $m$ is compatible with all the operations above.} To reduce a sum, a product, a substitution, or
an inverse, we may reduce the ingredients first and then compute in $\Z/m\Z$. This holds because every
coefficient of the result is obtained from finitely many coefficients of the ingredients by adding, subtracting and
multiplying, and these operations are compatible with remainders.

\para{Rings.} Some statements below hold over every commutative ring, for example $\Z$, $\Q$ and $\Z/m\Z$. No
division is used in them; this matters because in $\Z/m\Z$ division is not always possible (for instance, $2$ has no inverse in
$\Z/4\Z$).

\subsection{The Frobenius identity}
\begin{lemma}[Frobenius identity]\label{lem:frobenius}
Let $p$ be a prime. For every power series $U$ with coefficients in $\F_p$, we have $U(x)^p=U(x^p)$.
\end{lemma}

\begin{proof}
In $\F_p$, the binomial theorem simplifies to $(a+b)^p=a^p+b^p$, because every binomial coefficient
$\binom pi$ with $0<i<p$ is divisible by $p$. Also $c^p=c$ for every $c\in\F_p$ (Fermat's little theorem). So
$\bigl(\sum_nu_nx^n\bigr)^p=\sum_nu_n^px^{pn}=\sum_nu_nx^{pn}=U(x^p)$. The first equality holds for polynomials by
repeated use of $(a+b)^p=a^p+b^p$; it holds for power series because each coefficient of $U(x)^p$ involves only finitely
many coefficients of $U$.
\end{proof}

For $p=2$ this says that squaring a series modulo $2$ doubles every exponent: for instance,
$(x+x^3)^2\equiv x^2+x^6\pmod2$. In particular, the square of any series has, modulo $2$, no terms of odd degree.

\subsection{Doubling modulo 4}
\begin{lemma}\label{lem:double}
Let $U,V$ be series with coefficients in $\Z/4\Z$. Then $2U=2V$ if and only if $U$ and $V$ agree modulo $2$.
Consequently, if $P$ and $Q$ have constant term $0$ and $P\equiv Q\pmod 2$, then $2H(P)=2H(Q)$ in $\Z/4\Z$ for every
series $H$.
\end{lemma}

\begin{proof}
For residues $u,v$ modulo $4$, $2u=2v$ means that $4$ divides $2(u-v)$, that is, $2$ divides $u-v$. Apply this to
every coefficient. For the second claim, $H(P)$ and $H(Q)$ agree modulo $2$ because reduction is compatible with
substitution.
\end{proof}
In other words, \emph{twice a series, modulo $4$, depends only on the series modulo $2$.} This lets us use information modulo
$2$ to get information modulo $4$, without ever dividing by $2$ in $\Z/4\Z$, which is impossible.

\section{Results}\label{sec:results}

Let
\[
R(x)=\sum_{k\ge0}(-1)^kx^{2^k}=x-x^2+x^4-x^8+x^{16}-\cdots .
\]

\begin{theorem}[Exactly one solution]\label{thm:exist}
Equation \eqref{eq:oeis} has exactly one normalized solution $A$ with integer coefficients: the compositional inverse
of $R$. Its first coefficients are $a(1)=a(2)=1$.
\end{theorem}

\begin{theorem}[Hanna's conjecture]\label{thm:main}
For every $n>2$,
\[
a(n)\equiv\begin{cases}2\pmod4,&\text{if }n=3\cdot2^k\text{ for some }k\ge0,\\[2pt]0\pmod4,&\text{otherwise.}\end{cases}
\]
More precisely, $A\equiv x+x^2\pmod 2$, so that $A=x+x^2+2H$ for a series $H$ with integer coefficients, and
\[
H\equiv x^3+x^6+x^{12}+x^{24}+\cdots=\sum_{k\ge0}x^{3\cdot2^k}\pmod 2 .
\]
\end{theorem}

\section{Proofs}\label{sec:proofs}

\para{The idea in brief.} Formula (E5) of the entry turns the problem into a question about the inverse of the
series $R$. Modulo $2$, the equation $R(x)+R(x^2)=x$ shows that this inverse is $x+x^2$. To see one step further, we
write $A=x+x^2+2H$ exactly and compute modulo $4$. The doubling lemma (Lemma~\ref{lem:double}) lets us replace
complicated series inside $2H(\cdot)$ by their reductions modulo $2$, and \eqref{eq:oeis} collapses to
$H+H^2\equiv x^3\pmod 2$.

\begin{lemma}\label{lem:R}
$R$ is the only normalized series with $R(x)+R(x^2)=x$.
\end{lemma}

\begin{proof}
In $R(x)+R(x^2)=\sum_k(-1)^kx^{2^k}+\sum_k(-1)^kx^{2^{k+1}}$, every power $x^{2^j}$ with $j\ge1$ appears twice, with
opposite signs, so only $x$ remains. Conversely, in any normalized solution the coefficient of $x^n$, for $n\ge2$, must
be minus the coefficient of $x^{n/2}$ if $n$ is even, and $0$ if $n$ is odd. This determines all coefficients one at a
time, so the solution is unique.
\end{proof}

\begin{proof}[Proof of Theorem~\ref{thm:exist}]
\emph{Existence.} Let $A$ be the compositional inverse of $R$ (Lemma~\ref{lem:inverse}); it has integer coefficients.
Substitute $A(x)$ for $x$ in $R(x)+R(x^2)=x$. Since $R(A(x))=x$, this gives $x+R\bigl(A(x)^2\bigr)=A(x)$, that is,
$R\bigl(A(x)^2\bigr)=A(x)-x$. Substituting both sides into $A$ gives $A(x)^2=A\bigl(A(x)-x\bigr)$, which is
\eqref{eq:oeis}.

\emph{Uniqueness.} Let $U$ be any normalized integer solution of \eqref{eq:oeis}, and let $V$ be its inverse.
Substituting both sides of \eqref{eq:oeis} into $V$ gives $U(x)-x=V\bigl(U(x)^2\bigr)$. Substituting $V(x)$ for $x$
then gives $x-V(x)=V(x^2)$. So $V+V(x^2)=x$, hence $V=R$ by Lemma~\ref{lem:R}, and $U=A$.

\emph{First coefficients.} From $R=x-x^2+O(x^4)$ one computes $A=x+x^2+2x^3+O(x^4)$; for instance, the coefficient
of $x^2$ in $R(A(x))=x$ is $a(2)-1=0$.
\end{proof}

\begin{proof}[Proof of Theorem~\ref{thm:main}]
\emph{Step 1: modulo $2$.} By Lemma~\ref{lem:frobenius}, $R(x^2)\equiv R(x)^2\pmod 2$, so the equation of
Lemma~\ref{lem:R} becomes $R+R^2\equiv x$. This says that substituting $R$ into the polynomial $x+x^2$ gives $x$, so
$x+x^2$ is the compositional inverse of $R$ modulo $2$. By the uniqueness of inverses (Lemma~\ref{lem:inverse}),
$A\equiv x+x^2\pmod 2$. Hence there is a series $H$ with integer coefficients and no constant term such that
\[
A=x+x^2+2H\qquad\text{exactly.}
\]

\emph{Step 2: modulo $4$.} From now on we compute modulo $4$. Put $D=A-x=x^2+2H$. Then
\[
A(D)=D+D^2+2H(D),\qquad D^2=x^4+4x^2H+4H^2\equiv x^4 .
\]
Since $D\equiv x^2\pmod 2$, Lemma~\ref{lem:double} gives $2H(D)\equiv2H(x^2)$. So the left side of \eqref{eq:oeis} is
\[
A\bigl(A(x)-x\bigr)=A(D)\equiv x^2+2H+x^4+2H(x^2).
\]
The right side is $A^2=(x+x^2)^2+4(x+x^2)H+4H^2\equiv x^2+2x^3+x^4$. Comparing the two sides gives
\[
2\bigl(H+H(x^2)\bigr)\equiv2x^3\pmod 4 .
\]

\emph{Step 3: back to modulo $2$.} By Lemma~\ref{lem:double}, the last congruence means $H+H(x^2)\equiv x^3\pmod 2$,
and by Lemma~\ref{lem:frobenius}, $H(x^2)\equiv H^2$. So the reduction $\bar H$ of $H$ modulo $2$ satisfies
\[
\bar H+\bar H^2=x^3 .
\]
The series $S=\sum_{k\ge0}x^{3\cdot2^k}$ satisfies this equation: $S^2=\sum_{k\ge1}x^{3\cdot2^k}$, so all terms of
$S+S^2$ cancel in pairs except $x^3$. And there is only one solution without constant term, because the coefficient of
$x^n$ in $\bar H+\bar H^2$ is $\bar h_n$ plus $\bar h_{n/2}$ (when $n$ is even), which determines $\bar h_n$ from
earlier coefficients. Hence $\bar H=S$.

\emph{Step 4: conclusion.} For $n>2$ the polynomial $x+x^2$ contributes nothing to the coefficient of $x^n$, so
$a(n)=2\,\coeff{n}{H}$. This is $\equiv2\pmod 4$ when $\coeff{n}{H}$ is odd, that is, when $n=3\cdot2^k$, and
$\equiv0\pmod4$ otherwise.
\end{proof}

\section{Checking the results}\label{sec:checks}

A mathematical proof can contain mistakes, so the results were checked in two independent ways.

\para{Numerical check.}
We solved \eqref{eq:oeis} coefficient by coefficient with exact integer arithmetic, up to degree $129$. This
reproduces all 30 terms printed in the OEIS entry, and the rule of Theorem~\ref{thm:main} holds for every $n$ with
$2<n<130$.
The computation was done with the PARI/GP computer algebra system, using code that was written separately from the
formal proofs (file \texttt{verification/independent\_checks.gp} in the accompanying repository~\cite{repo}).
A computation can only test finitely many cases, so it is a check on the definitions and on the indexing, not a proof.

\para{Formal proof in Lean.}
Lean~\cite{lean} is a \emph{proof assistant}: a program in which mathematical statements and their proofs are written
in a precise formal language, and every step is checked mechanically down to the basic rules of logic. The theorems of Section~\ref{sec:results} are proved in Lean~4 (version 4.35.0) on top of
Mathlib~\cite{mathlib}, Lean's library of formalized mathematics. The statements are collected in the file \texttt{lean/Entries/A389537.lean} of the repository,
in the namespace \texttt{OEIS.A389537}. The full proofs are in \texttt{lean/Families/LacunaryInverse.lean}.
The table lists the results covered by the Lean formalization. The Lean formalization does not cover the additional arguments and observations in Section~\ref{sec:discussion}.
\begin{center}\small\renewcommand{\arraystretch}{1.25}
\begin{tabular}{>{\raggedright\arraybackslash}p{0.52\linewidth}>{\raggedright\arraybackslash}p{0.4\linewidth}}
\toprule
Statement in this paper & Lean theorem(s)\\
\midrule
Lemma~\ref{lem:R}: $R+R(x^2)=x$, and uniqueness & \texttt{auxiliary\_\allowbreak{}exists\_\allowbreak{}unique}, \texttt{auxiliary\_\allowbreak{}equation} \\
$A$ and $R$ are mutually inverse, and $A$ is normalized & \texttt{inverse\_\allowbreak{}right}, \texttt{inverse\_\allowbreak{}left}, \texttt{normalized} \\
Theorem~\ref{thm:exist}: $A$ solves \eqref{eq:oeis} and is the only normalized integer solution & \texttt{source\_\allowbreak{}equation}, \texttt{source\_\allowbreak{}exists\_\allowbreak{}unique}, \texttt{source\_\allowbreak{}unique} \\
Steps 1--3: $A\equiv x+x^2\pmod2$; $A=x+x^2+2H$; $\bar H+\bar H^2=x^3$; the odd coefficients of $H$ & \texttt{series\_\allowbreak{}mod\_\allowbreak{}two}, \texttt{integer\_\allowbreak{}decomposition}, \texttt{half\_\allowbreak{}equation\_\allowbreak{}mod\_\allowbreak{}two}, \texttt{half\_\allowbreak{}binary\_\allowbreak{}support} \\
Theorem~\ref{thm:main}: the remainders for $n>2$ & \texttt{residue\_\allowbreak{}mod\_\allowbreak{}four} \\
\bottomrule
\end{tabular}
\end{center}
The Lean project was rebuilt from source, and the axioms on which each of these theorems depends were listed with
Lean's \texttt{\#print axioms} command. Every one depends only on the three standard axioms of Lean's logic
(\texttt{propext}, \texttt{Classical.choice}, \texttt{Quot.sound}); no step is left unproved. A formal proof certifies
exactly the statement that Lean checks. Whether that statement says what the OEIS entry means is a separate question,
so the Lean definitions were also compared with the definition in the OEIS entry
(Section~\ref{sec:sequence}). The principal theorems are stated once more, using only Mathlib's definitions, in the
file \texttt{lean/Challenges/Halving.lean}, and Lean's Comparator tool confirms that the
proved theorems have exactly these statements.

\section{Discussion}\label{sec:discussion}

\para{Scope.} Theorem~\ref{thm:main} settles the conjecture completely, for every index. It concerns remainders modulo
$4$ only.

\para{The other observation in the entry.} The entry also observes that, for $n>7$, the signs of $a(n)$ come in runs
of length $9$ or $10$. This is a statement about the signs of the coefficients, not about their remainders, and our method
says nothing about it. Explaining it would require studying how fast the coefficients grow and how they oscillate, a
question of a different kind.

\para{Higher moduli.} One could try to continue the argument modulo $8$, by writing $H=x^3+x^6+\cdots+2H_2$ in the same
way. We have not done this.

\section*{About this paper}\label{sec:about}
\para{How the work was done.}
This paper is one of 28 papers on conjectures from the OEIS that came out of an AI research effort directed by
Dev~D.~Goyal~\cite{repo}. The roles were as follows.
\begin{itemize}
\item \emph{Dev~D.~Goyal} started the project, chose which problems to pursue, asked for proofs, formal verification and
revisions, and decided what to publish. Dev~D.~Goyal did not supply proof strategies or intermediate arguments.
\item \emph{ChatGPT} (OpenAI), using the GPT-6 model with Pro reasoning, worked autonomously from those prompts. It
searched the OEIS for suitable conjectures, found the proofs, wrote the Lean~4 formalization, and wrote a first draft
of this paper.
\item \emph{Claude Opus~5.5}, used with Claude Code, checked the work independently. It rebuilt the Lean project from source and confirmed that the theorems listed in Section~\ref{sec:checks} rely only on
Lean's standard axioms; compared the Lean definitions with the OEIS definition; recomputed the sequence with separately
written PARI/GP code, compared it with the terms published in the OEIS, and tested every claim of the paper on many
terms; and read every step of the proofs. It then wrote the present version of the paper, expanding the
explanations and, in places, writing out an argument in more detail or by a more elementary route. The theorems of
Section~\ref{sec:results} and the proof strategies are ChatGPT's. The Discussion section, including the short arguments
in it, was written by Claude Opus~5.5; its observations were checked by computer and are not part of the Lean
formalization.
\end{itemize}
Dev~D.~Goyal has reviewed the proofs and is responsible for this paper; this review is not an independent expert
validation.

\para{Review status.} This paper has not been peer reviewed, and it has not been submitted to a journal or to the OEIS.

\para{Priority.} An OEIS entry that labels a statement as a conjecture shows only that its contributor did not know
a proof; it does not show that no proof existed. During the research only limited searches for earlier work were
made. No claim of priority is made.

\para{Corrections.} Please report any error in the issue tracker of the repository~\cite{repo}, where the Lean sources and the numerical checks are also
available.

\section*{Acknowledgements}
The sequence, the conjecture, and formula (E5) identifying $A$ as the inverse of the lacunary series $R$ are due to
Paul~D.~Hanna~\cite{oeis}.

\begin{thebibliography}{99}\small
\raggedright
\bibitem{oeis} P.~D.~Hanna, entry A389537 (G.f.\ $A(x)$ satisfies $A(A(x)-x)=A(x)^2$), \emph{The On-Line Encyclopedia of Integer Sequences}, November 2025. \url{https://oeis.org/A389537}. Text as of October 2026.
\bibitem{allouche} J.-P.~Allouche and J.~Shallit, \emph{Automatic Sequences}, Cambridge University Press, 2003.
\bibitem{christol} G.~Christol, T.~Kamae, M.~Mend\`es France and G.~Rauzy, Suites alg\'ebriques, automates et substitutions, \emph{Bulletin de la Soci\'et\'e Math\'ematique de France} 108 (1980), 401--419.
\bibitem{oeisdb} N.~J.~A.~Sloane and The OEIS Foundation Inc., \emph{The On-Line Encyclopedia of Integer Sequences}. \url{https://oeis.org}.
\bibitem{stanley} R.~P.~Stanley, \emph{Enumerative Combinatorics}, Volume~2, Cambridge University Press, 1999. (Chapter~5 treats formal power series, composition and compositional inverses.)
\bibitem{niven} I.~Niven, Formal power series, \emph{American Mathematical Monthly} 76 (1969), 871--889.
\bibitem{lean} L.~de~Moura and S.~Ullrich, The Lean~4 theorem prover and programming language, in \emph{Automated Deduction -- CADE 28}, Lecture Notes in Computer Science 12699, Springer, 2021, 625--635.
\bibitem{mathlib} The mathlib Community, The Lean mathematical library, in \emph{Proceedings of the 9th ACM SIGPLAN International Conference on Certified Programs and Proofs (CPP 2020)}, 367--381. Version used: commit \texttt{56a005855ea1}.
\bibitem{repo} D.~D.~Goyal, \emph{Proofs and corrections of OEIS conjectures}, repository, 2026. \url{https://github.com/dvgyl/oeis-proofs}.
\end{thebibliography}
\end{document}
